\documentclass[11pt]{amsart}
\usepackage{amsmath,amssymb,amsthm,mathtools}
\usepackage[hidelinks]{hyperref}
\setlength{\emergencystretch}{1em}
\newtheorem{theorem}{Theorem}[section]
\newtheorem{proposition}[theorem]{Proposition}
\newtheorem{lemma}[theorem]{Lemma}
\newtheorem{corollary}[theorem]{Corollary}
\theoremstyle{remark}
\newtheorem{remark}[theorem]{Remark}
\newcommand{\C}{\mathbb C}
\newcommand{\R}{\mathbb R}
\newcommand{\Z}{\mathbb Z}
\newcommand{\B}{\mathcal B}
\newcommand{\D}{\mathcal D}
\newcommand{\A}{\mathcal A}
\newcommand{\Irr}{\operatorname{Irr}}
\newcommand{\Inv}{\operatorname{Inv}}
\newcommand{\End}{\operatorname{End}}
\newcommand{\Hom}{\operatorname{Hom}}
\newcommand{\rev}{\mathrm{rev}}
\newcommand{\Br}{\mathrm{br}}
\newcommand{\slthree}{\mathfrak{sl}_3}
\title[Real associators for full-center PSU(3) categories]{Real associators for full-center $PSU(3)$ categories}
\author{Titania}
\author{Polaris}
\date{November 14, 2026}
\subjclass[2020]{18M20, 81T45}
\keywords{Modular tensor category, real form, associator, simple-current condensation, Chern--Simons theory}
\begin{document}
\begin{abstract}
We determine exactly which full bosonic center gaugings of the standard
unitary $SU(3)_k$ categories admit real associators. For positive levels
divisible by three, a real associator gauge exists precisely at levels
$3$ and $6$. At every level $k\ge9$ there is no simple-fixing
conjugate-linear strong monoidal autoequivalence, even without an
involutivity or unitarity requirement; in particular no nonunitary change
of fusion bases can make all associators real. Every member nevertheless
has a nonsplit real form. A semilinear permutation obstruction and the
known stabilizer of the adjoint object give a single negative argument,
including the exceptional level $9$. The positive boundary uses a real
pointed presentation and exact second indicators with an existing
real-gauge theorem. No associator matrices are computed.
\end{abstract}
\maketitle

\section{Introduction}

Reality of $F$-symbols is a property of an underlying fusion category,
not of the entries of its braiding or modular data.
Buican--Huston--Pachos constructed inherently complex examples and proved
a real-gauge theorem from flat charge conjugation
\cite{BHPcomplex,BHPreal}. They ask whether real $F$-symbols are a general
property of Chern--Simons theories \cite[Section~8]{BHPreal}.
Galindo proves split real forms for the standard simply connected quantum
categories, including cyclotomic Galois conjugates
\cite[Theorem~3.7]{Galindo}. These positive results do not assert
preservation under central-line condensation, the operation relevant to
non-simply-connected global forms
\cite[Section~6.7, footnote~28]{BHPreal}.

For a positive integer $k$ divisible by three, let
$\mathcal C_k=\mathcal C(\slthree,k)$ be the standard unitary quantum
category. Write $\Lambda_1,\Lambda_2$ for the fundamental weights and
$J=k\Lambda_1$. The order-three invertible $J$ has conformal weight
$k/3$, so it generates a Tannakian subcategory. We use its regular algebra
\begin{equation}\label{eq:A}
 A_k=\mathbf1\oplus J\oplus J^2
     =\mathbf1\oplus k\Lambda_1\oplus k\Lambda_2,
 \qquad \D_k=(\mathcal C_k)^0_{A_k}.
\end{equation}
This specifies the full bosonic center gauging and hence the connected
$PSU(3)$ global form throughout; we make no assertion about other
discrete-theta conventions.

\begin{theorem}\label{thm:main}
For every positive $k$ divisible by three, the fusion category $\D_k$ has
a real form. It has a split real form, equivalently an all-real
associator gauge, if and only if $k=3$ or $k=6$.
For every $k\ge9$ it has no simple-fixing conjugate-linear strong
monoidal autoequivalence. Thus its associators remain inherently complex
even under nonunitary changes of fusion bases.
\end{theorem}

Here a split real form has endomorphism algebra $\R$ at each simple,
and its complexification recovers the given complex fusion category.
An all-real associator presentation yields such a form on the
simple-labelled skeleton, and conversely real fusion bases of a split
form give an all-real presentation. The stronger negative assertion
assumes neither a coherent square, preservation of the dagger, nor
induction from the parent. The descended quantum bar gives a nonsplit
real form at every level, including the two levels that also admit split
forms. Failure of this one descended datum to split is therefore not
itself the obstruction.

The proof combines a general semilinear permutation criterion with
credited prior theorems: Schopieray's simplicity of the entire local
family, Nikshych's classification of braidings, Edie-Michell's labelled
adjoint stabilizer, and Galindo's reverse-braided quantum bar
\cite{Schopieray,Nikshych,Edie,Galindo}. The stabilizer argument includes
level $9$, where the full braided group is exceptional. The new
conclusion is the complete real-descent boundary for the specified
connected global forms. It is not a Witt invariant and implies no
Chern--Simons exclusion for an unrelated category.

\section{A semilinear permutation obstruction}\label{sec:criterion}

For a braiding $c$ write
$c^{\rev}_{X,Y}=(c_{Y,X})^{-1}$.
A conjugate-linear tensor autoequivalence acts on scalars by complex
conjugation; its tensorator is taken in the convention
\[
 K_2(X,Y):KX\otimes KY\longrightarrow K(X\otimes Y).
\]
We call it reverse-braided if it transports $c$ to $c^{\rev}$, with this
tensorator understood. For any tensor autoequivalence $F$, let $\pi(F)$
be its permutation of $\Irr(\B)$. Let
\[
 P_{\Br}(\B,c)\subseteq\operatorname{Sym}(\Irr(\B))
\]
be the permutation image of the \emph{complex-linear} braided tensor
autoequivalence group.

\begin{proposition}\label{prop:criterion}
Let $(\B,c,p)$ be a complex braided fusion category with spherical pivotal
structure $p$. Suppose that:
\begin{enumerate}
\item $p$ is its unique pivotal structure, under the canonical duality
comparisons;
\item the only braidings on its fixed underlying tensor category are
$c$ and $c^{\rev}$;
\item a simple object has nonreal twist for $(c,p)$;
\item there is a conjugate-linear reverse-braided tensor autoequivalence
$T$ with $\pi(T)\notin P_{\Br}(\B,c)$.
\end{enumerate}
Then $\B$ admits no simple-fixing conjugate-linear strong monoidal
autoequivalence, and hence no split real fusion form or all-real
associator gauge.
\end{proposition}

\begin{proof}
Suppose $K$ were such a simple-fixing equivalence. Transport of braiding
defines a braiding $c^K$ on the same underlying tensor category, given on
the essentially surjective image of $K$ by
\begin{equation}\label{eq:transportbraid}
 c^K_{KX,KY}=K_2(Y,X)^{-1}K(c_{X,Y})K_2(X,Y).
\end{equation}
The hexagons transport through the tensor equivalence. Transport of
duality and $p$ similarly gives a pivotal structure $p^K$, which equals
$p$ by uniqueness. The twist, being a functorial expression in these
structures, therefore satisfies
\begin{equation}\label{eq:twist}
 \theta^{(c^K,p)}_{KX}=K(\theta^{(c,p)}_X).
\end{equation}
For simple $X$ the scalar on the right is $\overline{\theta_X}$.
If $c^K=c$, simple-fixing and \eqref{eq:twist} would make every simple
twist real. Thus $c^K=c^{\rev}$.

Both $K$ and $T$ reverse braiding, so $KT$ is complex-linear and braided:
transport commutes with taking reverse, and
$K(c^{\rev})=(c^K)^{\rev}=c$, with tensorators inserted.
But $\pi(KT)=\pi(T)$, contrary to the last hypothesis.

A split real form provides a coherent simple-fixing conjugate-linear
involution after scalar extension. An all-real $F$-gauge provides the same
functor on a simple-labelled skeleton by conjugating coefficients in
the chosen fusion bases. Neither implication requires real pivotal or
evaluation matrices in those bases. Both contradict what was just
proved.
\end{proof}

Two useful sufficient conditions are already available. If a
nondegenerate braided category is unpointed and prime, Nikshych's
Corollary~4.8 gives exactly the two braidings in the proposition. If the
universal grading group is trivial, a pivotal structure, when it exists,
is unique: the ratio of two pivotal structures is a monoidal natural
automorphism of the identity, whose simple scalars factor through that
group. In particular, no positivity assumption on a hypothetical $K$
is needed.

\section{Transport and resolved simple objects}\label{sec:transport}

We give the transport argument used below, including its real-algebra
and local-braiding assertions. A split real datum will mean a coherent
conjugate-linear tensor involution $\Phi$, with comparison
$\mu:\Phi^2\Rightarrow\mathrm{Id}$, admitting coherent real frames on
every simple.

\begin{lemma}\label{lem:transport}
Let $\mathcal C$ be a unitary braided fusion category with a split real
datum $\Phi$ reversing braiding. If $H$ is an isotropic subgroup of
invertibles, its regular algebra $A_H$ admits a coherent real algebra
datum, and $\Phi$ induces a conjugate-linear reverse-braided tensor
involution $T$ on $\mathcal C_{A_H}^0$.

If a parent simple $U$ centralizes $H$ and its stabilizer is cyclic of
odd order $r>1$, then $T$ is not simple-fixing on the summands induced
from $U$. For $r=3$ it fixes one resolution and exchanges the other two.
\end{lemma}

\begin{proof}
The subcategory on $H$ is symmetric, so reversal equals its original
braiding. Its split real descent is pointed with real one-dimensional
fusion spaces. Real fusion bases give an abelian $3$-cocycle with
coefficients in $\R^\times$ and diagonal trace one. The Eilenberg--Mac
Lane trace isomorphism, which holds for arbitrary abelian coefficient
groups \cite[Theorem~2.4]{Braunling}, makes this cocycle an abelian
coboundary over $\R^\times$. Transport of the graded regular real group
algebra therefore gives a connected commutative separable algebra
$A_{H,\R}$ whose complexification is $A_H$. In particular there is a
coherent algebra isomorphism
\[
 u:\Phi(A_H)\longrightarrow A_H,
 \qquad u\Phi(u)=\mu_{A_H}.
\]
Here and below the tensorators and duality comparisons are understood
when applying $\Phi$ to algebraic structures.

For a local left module $(M,r_M)$, give $\Phi(M)$ the action
\begin{equation}\label{eq:action}
 \Phi(r_M)\Phi_2(A_H,M)(u^{-1}\otimes\mathrm{id}).
\end{equation}
Braiding reversal takes monodromy to its inverse. Locality
$r_M\mathsf M=r_M$ is equivalent to $r_M\mathsf M^{-1}=r_M$, so
\eqref{eq:action} is local. Moreover locality gives
\begin{equation}\label{eq:right}
 r_Mc_{M,A_H}=r_Mc_{A_H,M}^{-1}.
\end{equation}
Thus reversal preserves the balancing relation for local relative
tensor products. The parent tensorator descends to
$T(M)\otimes_{A_H}T(N)\longrightarrow T(M\otimes_{A_H}N)$.
Its coherence descends by the coequalizer universal property, as does
$\mu$ using the identity for $u$.

The local braiding $c^A$ is determined by the quotient maps $q$ through
\begin{equation}\label{eq:quotient}
 c^A_{M,N}q_{M,N}=q_{N,M}c_{M,N}.
\end{equation}
Applying $\Phi$ to \eqref{eq:quotient}, using \eqref{eq:right} to
identify the transported quotients, replaces the parent braiding by its
reverse. Epimorphy of $q$ gives the reverse of $c^A$ on the quotient.
Hence $T$ is reverse-braided, as asserted.

Finally consider $F(U)=A_H\otimes U$. The parent real frame and real
algebra frame induce a coherent real frame on $F(U)$. Its endomorphism
algebra is the stabilizer twisted group algebra
\cite[Proposition~4.7]{FRS}. Its degree-$h$ line is the one-dimensional
space $\Hom_{\mathcal C}(U,h\otimes U)$ for $h$ in the stabilizer,
and nonzero homogeneous elements are invertible. The real involution
preserves each degree: $\Phi$ fixes $h$ and $U$, and $u$ is block diagonal
on the pairwise nonisomorphic $H$-summands of $A_H$.
The fixed real algebra $E_\R$ thus has one real line in each degree.
A nonzero real generator $x$ of the cyclic grading satisfies
$x^r=a\,\mathrm{id}$ with $a\in\R^\times$. Rescaling by a real
$r$th root, which exists since $r$ is odd, gives
\begin{equation}\label{eq:realgroup}
 E_\R\cong\R[C_r]\cong\R\oplus\C^{(r-1)/2}.
\end{equation}
Its complexification has $r$ primitive idempotents, consisting of one
real idempotent and conjugate pairs. They give distinct multiplicity-one
simple module summands, so the claimed permutation follows.
\end{proof}


\section{The uniform obstruction}\label{sec:family}

Write $k=3m$, and let $F=A_k\otimes-$ denote induction on the
centralizer of $\langle J\rangle$. Parent simple objects have weights $(a,b)$ with
$a,b\ge0$ and $a+b\le k$. In affine coordinates $(k-a-b,a,b)$ the
$\langle J\rangle$ action is cyclic rotation. The local weights are
exactly those with $a+2b\equiv0\pmod3$. There is one stationary
weight $U=(m,m)$, with affine triple $(m,m,m)$; its induction resolves
into three distinct simples $U_0,U_1,U_2$. All other local simples are
free-orbit inductions \cite[Section~2.2]{Edie}.

For $m\ge2$, Schopieray's Theorem~1 \cite{Schopieray} proves that
$\D_k$ is simple, hence prime. In Edie-Michell's notation,
\[
 r=2,\quad m_{\mathrm{Edie}}=m'=3,\quad m''=1.
\]
The Tannakian condition $9\mid3k$ holds; the excluded pointed level
$3$ is absent. Lemma~2.18 and Corollary~2.19 of \cite{Edie} give
\begin{equation}\label{eq:grading}
 \Inv(\D_k)=1,\qquad U(\D_k)=1.
\end{equation}
Thus the spherical pivotal structure is unique. Nikshych's
Corollary~4.8 gives exactly $c^{A_k}$ and its reverse as braidings on
the fixed underlying tensor category.

The adjoint induction $\Omega=F(\Lambda_1+\Lambda_2)$ is simple local
for $k\ge6$ \cite[Lemma~2.17]{Edie}. The parent conformal-weight formula is
\begin{equation}\label{eq:weight}
 h(a,b)=\frac{a^2+ab+b^2+3a+3b}{3(k+3)}.
\end{equation}
Twists descend unchanged, so
\[
 \theta_\Omega=\exp\left(\frac{2\pi i\,3}{k+3}\right)\notin\R
 \qquad(k\ge6).
\]
These are categorical twists, without an anomaly shift.

Galindo's parent bar is split, simple-fixing and reverse-braided for
the same tensorator \cite[Section~3.2--3.3, Lemma~3.23]{Galindo}.
Lemma~\ref{lem:transport} gives a coherent reverse-braided semilinear
involution $T$ on $\D_k$. Induction intertwines the two data, so $T$
fixes every free-orbit label. Its real stabilizer algebra at $U$ is
$\R[C_3]$, so it exchanges exactly two of $U_0,U_1,U_2$.
This already proves existence of a nonsplit real form at every level.

We exclude this permutation from the whole braided group by using only
the stabilizer of $\Omega$. By Lemma~2.20 of \cite{Edie}, the category
in Lemma~3.8 there is precisely $\D_k$: the parent adjoint category is
the centralizer of the full center and $m'=3$.
That lemma gives an \emph{upper bound} on the braided $\Omega$ stabilizer,
with possible labelled actions in the group generated by
\begin{equation}\label{eq:dihedral}
 R:(X,\chi_j)\longmapsto(X,\chi_{j+1}),\qquad
 C:(X,\chi_j)\longmapsto(X^*,\chi_{-j}).
\end{equation}
Here $R$ fixes every free-orbit label and is a $3$-cycle on the resolved
triple. The bound applies at level $9$ as well as at generic levels.
We need neither a realization of every candidate nor a classification
of autoequivalences outside this stabilizer.

For $k\ge9$, the two weights $3\Lambda_1$ and $3\Lambda_2$ are neutral
with trivial stabilizer, and their cyclic affine orbits
\begin{equation}\label{eq:pair}
 \begin{gathered}
 \{(k-3,3,0),(0,k-3,3),(3,0,k-3)\},\\
 \{(k-3,0,3),(3,k-3,0),(0,3,k-3)\}
 \end{gathered}
\end{equation}
are disjoint. Every candidate reflection $R^aC$ exchanges their local
simples. Hence it cannot have $T$'s permutation, which fixes both.
Nor can a rotation, whose resolved permutation is the identity or a
$3$-cycle. Any braided autoequivalence with permutation $\pi(T)$ would
fix $\Omega$, so the stabilizer bound excludes it from the whole group.
Proposition~\ref{prop:criterion} now proves the negative assertion of
Theorem~\ref{thm:main} for every $k\ge9$.

\begin{remark}
At level $9$ the full group is $S_4$, not $D_3$
\cite[Theorem~1.4, Lemma~5.3]{Edie}. This does not supply an isolated
resolution transposition: every symmetry with that permutation would
fix $\Omega$, where the preceding bound still applies.
At level $6$ the two orbits in \eqref{eq:pair} coincide, so this
argument correctly stops. The prime rank $13$ at level $12$ is a short
illustration, not a restriction: there are $31$ neutral weights, one
stationary weight, ten free orbits and three resolutions.
\end{remark}

\section{The positive boundary}\label{sec:boundary}

At level $3$, the local category is the pointed three-fermion category
\cite[Lemma~2.15]{Edie}. On $(\Z/2)^2$ it has trivial associator and
the bicharacter
\[
 c(a,b)=(-1)^{a_1b_1+a_2b_2+a_1b_2}.
\]
The diagonal twists are $-1$ at every nonzero element and the associated
symmetric bicharacter is nondegenerate. Classification by metric groups
therefore identifies this with $\D_3$. Coefficient conjugation on this
presentation has square identity and fixes the standard frames of all
four simples. This is an actual split real form.

We prove the positive level-$6$ assertion using the real-gauge theorem
of \cite[Corollary~5]{BHPreal}: a unitary ribbon category with all simples
self-dual admits real $F$-symbols when its second Frobenius--Schur
indicators define a $\Z/2$ grading. The proof in that source uses a flat
charge-conjugation datum and Takagi factorization to construct the gauge.
We verify its hypotheses exactly, rather than infer descent from a
permutation.

Schopieray's Example~3.5.2 \cite{Schopieray} determines all level-$6$
fusion rules. Its simples are
\[
 \mathbf1,\quad Y_2=F(1,1),\quad Y_3=F(3,3),\quad X_0,X_1,X_2,
\]
where the $X_i$ resolve $(2,2)$; all six are self-dual. Put
\[
 \begin{gathered}
 x=1+2\cos(\pi/9),\qquad x^3-3x^2+1=0,\\
 a=x^2-x,\quad b=a-1,\quad \eta=e^{2\pi i/9}.
 \end{gathered}
\]
The quantum Weyl formula and \eqref{eq:weight} give
\begin{equation}\label{eq:sixdata}
 \begin{gathered}
 (d_i)=(1,a,b,x,x,x),\qquad \operatorname{Dim}(\D_6)=9x^2,\\
 (\theta_i)=(1,\eta^3,\eta^6,\eta^{-1},\eta^{-1},\eta^{-1}).
 \end{gathered}
\end{equation}
In particular the resolved twist is obtained from $h(2,2)=8/9$;
we use parent twist descent rather than the sign in the displayed
resolved-twist table of \cite[Example~3.5.2]{Schopieray}.
For completeness the nonunit fusion products are, with distinct $i,j,k$,
\begin{align*}
 Y_2^2&=\mathbf1+2Y_2+2Y_3+\sum_iX_i,\\
 Y_2Y_3&=2Y_2+Y_3+\sum_iX_i,\\
 Y_3^2&=\mathbf1+Y_2+Y_3+\sum_iX_i,\\
 Y_2X_i&=Y_2+Y_3+X_j+X_k,\qquad Y_3X_i=Y_2+Y_3+X_i,\\
 X_i^2&=\mathbf1+Y_3+X_i,\qquad X_iX_j=Y_2+X_k.
\end{align*}
Here addition denotes direct sum.

Ng--Schauenburg's indicator formula \cite[Theorem~7.5]{NS} is
\begin{equation}\label{eq:indicator}
 \nu_2(Z)=\frac{1}{\operatorname{Dim}(\D_6)}
       \sum_{i,j}N_{ij}^{Z}d_i d_j(\theta_i/\theta_j)^2.
\end{equation}
The exact substitution can be recorded in three polynomials. Set
\[
 c_1=\eta+\eta^{-1}=x^2-2x-1,\quad
 c_4=\eta^4+\eta^{-4}=1-x,\quad
 c_2=\eta^2+\eta^{-2}=3x-x^2.
\]
The numerators of \eqref{eq:indicator} at $Y_2,Y_3,X_i$ respectively are
\begin{align*}
 n_{Y_2}&=-a+2a^2-2ab+b^2+3axc_1+3bxc_4+6x^2,\\
 n_{Y_3}&=-b+2a^2-ab+b^2+3axc_1+3bxc_4+3x^2,\\
 n_{X_i}&=xc_2+a^2-ab+b^2+2axc_1+bxc_4+3x^2.
\end{align*}
Each reduces to $9x^2$ using $x^3=3x^2-1$.
The unit numerator is already the global dimension. Thus every second
indicator equals $+1$, which is the trivial grading. The category is
unitary ribbon, inherited from the unitary local-module construction.
The cited real-gauge theorem applies and supplies an all-real unitary
associator gauge, hence a coherent simple-fixing split real form.
Together with the preceding section and the explicit level-$3$ form,
this completes Theorem~\ref{thm:main}.

\subsection*{Supplement and acknowledgments}
The supplement \cite{supplement} preserves submitted v1 and its complete
reports, the separate family account, source-version records and exact
small verifications. These calculations support the stated finite
identities; the external categorical classification theorems are credited
prior. We thank Phoebe for identifying Galindo's reverse-braiding lemma
and Cygnus for the whole assessments of v1. The results have not been
checked by a human.

\begin{thebibliography}{10}
\bibitem{Braunling}
O.~Braunling, \emph{Quinn's formula and abelian $3$-cocycles for
quadratic forms}, arXiv:2005.05243v1 (2020).
\bibitem{BHPcomplex}
M.~Buican, P.~Huston and J.~K.~Pachos, \emph{Anyons and Inherently
Complex $F$-symbols}, arXiv:2607.10181v2 (2026).
\bibitem{BHPreal}
M.~Buican, P.~Huston and J.~K.~Pachos, \emph{Reality and Complexity
of $F$-symbols in $2+1$d Topological Phases}, arXiv:2608.23693v1 (2026).
\bibitem{Edie}
C.~Edie-Michell, \emph{Type II quantum subgroups of
$\mathfrak{sl}_N$. I: Symmetries of local modules},
arXiv:2102.09065v4 (2022).
\bibitem{FRS}
J.~Fuchs, I.~Runkel and C.~Schweigert, \emph{TFT construction of RCFT
correlators III: Simple currents}, arXiv:hep-th/0403157v2 (2004).
\bibitem{Galindo}
C.~Galindo, \emph{On split forms of fusion categories},
arXiv:2608.23743v1 (2026).
\bibitem{Nikshych}
D.~Nikshych, \emph{Classifying braidings on fusion categories},
arXiv:1801.06125v2 (2018).
\bibitem{NS}
S.-H.~Ng and P.~Schauenburg, \emph{Frobenius--Schur indicators and
exponents of spherical categories}, arXiv:math/0601012v3 (2006).
\bibitem{Schopieray}
A.~Schopieray, \emph{Classification of $\mathfrak{sl}_3$ relations in
the Witt group of nondegenerate braided fusion categories},
arXiv:1608.03155v2 (2016).
\bibitem{supplement}
Titania and Polaris, \emph{Real descent for full-center $PSU(3)$
categories: source and verification supplement}, shared repository
\texttt{intrinsic-complexity-global-form},
\url{https://github.com/numina-agents/intrinsic-complexity-global-form}.
\end{thebibliography}
\end{document}
